\chapter{Federated Learning}
In questa sezione viene introdotto il Federated Learning, come funziona,
i diversi possibili approcci, quali sono i vantaggi e gli svantaggi.
Prima di tutto viene spiegato cos'è il Federated Learning e i diversi tipi
di partizionamento dei dati, mentre nel secondo paragrafo vengono discussi
gli aspetti positivi di questa tecnica, specie nell'area della privacy dei
dei dati per contesti ad alta sensibilità. Infine vengono discusse le
difficoltà che quest'approccio provoca, principalmente la distribuzione non
IID (Indipendent and Identically Distributed) dei dataset locali.

\section{Come funziona}
Il Machine Learning è la classe di algoritmi di apprendimento in cui ad un
modello viene fornito un insieme di dati, il training dataset, da cui
imparare una certa funzione obbiettivo su questi dati. Un algoritmo di 
Machine Learning popolare una decina di anni fa era quello delle SVM
(Support Vector Machine), mentre ad oggi il Deep Learning la fa da padrona.
Tipicamente, nei problemi di Machine Learning si assume di avere un unico
dataset e un unico modello che ha accesso a tutti i sample nel dataset.
Tuttavia quest'approccio può porre alcuni problemi: in contesti in cui si
lavora con dati sensibili, come i dati medici di privati cittadini, può 
essere sconsigliato, difficile o vietato (da norme come il GDPR o il più
recente AI act) raccogliere tali dati da utilizzare per allenare un modello;
inoltre la creazione e gestione di un tale dataset può avere costi non
indifferenti dati dal trasporto dei dati in unico data storage centralizzato,
la memorizzazione di tali dati e il costo dell'effettivo allenamento del
modello.

Il Federated Learning si pone di risolvere questo problema allenando il nostro
modello proprio alla fonte dei dati, dove questi sono già presenti o vengono
generati. In questo framework possiamo individuare N client ognuno dei quali
ha una copia del modello da allenare fornitagli da un server e ognuno dei quali
con il proprio dataset, disgiunto da quello degli altri client.
Ad ogni round di apprendimento ognuno degli N client allena il proprio modello 
con il proprio dataset locale e, finito il passo di apprendimento, invia al
server il proprio modello aggiornato. A questo punto il server si occuperà 
semplicemente di fare una qualche aggregazione dei nuovi N modelli che cli sono
stati forniti e redistribuire ai client il nuovo modello prodotto.

In questo studio ci occupiamo di applicare il Federated Learning a 2 reti neurali,
un semplice MLP (Multi Layer Perceptron) e una CNN (Convolutional Neural Network),
ma questa tecnica può essere usata in tutti gli algoritmi di Machine Learning il
cui modello usato permette un qualche meccanismo di aggregazione; le reti neurali
non sono uniche in questo e le già citate SVM sono un altro di questi casi.

A seconda di come differiscono i dataset locali di ogni client si possono
fare 2 classificazioni di Federated Learning: quello a partizionamento orizzontale
e quello a partizionamento verticale.

\section{Partizionamento orizzontale}
Nel partizionamento orizzontale i dataset di tutti i client hanno lo
stesso schema, le stesse features.

\section{Partizionamento verticale}
Nel partizionamento verticale invece ogni client utilizza un insieme 
di feature diverse da quelle degli altri, con un qualche collegamento 
tra i sample.

\section{Ibridazione}

\section{Vantaggi}

\section{Problemi}

Here is a citation~\cite{einstein1905}.
